\chapter*{Over het Open Logic Project}
\addcontentsline{toc}{chapter}{Over het Open Logic Project}


De \textit{Open Logic Text} is een opensourceleerboek waaraan mensen samen
schrijven. Het gaat over formele
metalogica, dus over het onderzoek naar formele logische systemen, en over
formele methoden. Het boek begint niet bij nul:
het gaat ervan uit dat je al een inleidende cursus formele logica hebt gevolgd.
Het boek is bedoeld voor lezers buiten de wiskunde, met name voor studenten
filosofie en informatica. Toch behandelt het de stof wiskundig precies.

Sommige onderwerpen staan al in het boek, maar zijn misschien nog niet
volledig behandeld. Veel stukken moeten nog flink worden herzien. We zijn van
plan het boek later met meer onderwerpen uit te breiden. Ook zijn we van plan
onderdelen aan het boek toe te voegen, bijvoorbeeld een woordenlijst, een lijst
met tips voor wat je verder kunt lezen, uitleg over de geschiedenis,
afbeeldingen en betere uitleg. Verder zijn we van plan stukken toe te voegen
waarin we uitleggen waarom de resultaten belangrijk zijn voor filosofie,
informatica en wiskunde, en meer opgaven en voorbeelden.
Zie je een fout of heb je een voorstel,
\href{https://github.com/OpenLogicProject/OpenLogic/wiki/Contributing}{laat het projectteam dat dan weten}.

Het project werkt volgens het idee van open source. De tekst is niet alleen
vrij beschikbaar; we geven ook de LaTeX-broncode vrij onder de licentie
Creative Commons Naamsvermelding. Volgens die licentie mag iedereen ons werk
downloaden, gebruiken, aanpassen, anders indelen, omzetten en verder
verspreiden, zolang de makers op de juiste manier worden genoemd. Meer
informatie staat op de website van het
Open Logic Project: \href{http://openlogicproject.org/}{openlogicproject.org}.
